\chapter{Considerações Finais}

\section{Conclusão}


Ao longo deste trabalho, foi possível aplicar conceitos fundamentais da Engenharia de Software, incluindo a definição de requisitos, modelagem de sistemas e análise de metodologias de desenvolvimento. Conforme descrito no Capítulo~\ref{cap:introducao}, o objetivo principal do projeto é a criação de um aplicativo móvel voltado à promoção da saúde mental e hábitos saúdaveis no contexto universitário. No entanto, a arquitetura desenvolvida apresenta flexibilidade suficiente para ser adaptada a outras plataformas, como aplicações web ou desktop. Isso é possível devido à estrutura robusta e desacoplada do \textit{backend}, que foi projetado para garantir escalabilidade e portabilidade.

Adicionalmente, a pesquisa bibliográfica realizada, aliada à orientação de uma coorientadora com formação em Psicologia, contribuiu significativamente para o aprofundamento do conhecimento sobre saúde mental. Essa parceria interdisciplinar foi essencial para compreender as necessidades do público-alvo e delinear funcionalidades que realmente dialoguem com os desafios enfrentados pelos estudantes universitários. Assim, o projeto não apenas se beneficia de uma base técnica sólida, como também de uma fundamentação teórica consistente na área da saúde.


\section{Trabalhos Futuros}

Diante dos resultados alcançados nesta primeira etapa, torna-se viável o desenvolvimento da segunda etapa (TCC2), na qual serão implementadas as funcionalidades definidas com base nos requisitos funcionais e não funcionais levantados nesta fase inicial.

Em etapas futuras, espera-se que o projeto evolua para a construção completa do aplicativo proposto, com a implementação das funcionalidades detalhadas no protótipo de alta fidelidade apresentado no Apêndice~\ref{apendice:prototipo}. Além disso, é recomendável uma vasta gama de usuários, envolvendo diferentes perfis da comunidade universitária. Isso permitirá conduzir testes de usabilidade, desempenho e validação com maior rigor metodológico.

Outra possibilidade relevante consiste na integração do sistema com soluções já existentes na área da saúde digital. Um exemplo interessante seria a interoperabilidade com o aplicativo \textit{Guardiões da Saúde}, utilizado pela Faculdade de Ciências da Saúde da UnB, que adota uma abordagem de vigilância participativa. A incorporação de funcionalidades complementares, como alertas de saúde, registros de hábitos e indicadores personalizados, poderia potencializar o alcance da proposta.

% FUI DORMIR É ISSO :)
